\begin{abstract}
\noindent\textbf{Purpose:} This is a \emph{pre-analysis study protocol} for STARS (Stress, Tinnitus, and Audiometric Research Study), not an empirical report. It prespecifies public-data analyses of the associations among \emph{perceived stress and work-related factors}, tinnitus, and hearing outcomes in adults, together with a cross-national external-validation plan and a clinician-governed, safety-gated artificial-intelligence (AI) architecture. We deliberately avoid causal overclaiming: because the population surveys do not contain validated occupational-stress instruments, the exposure is labeled perceived stress and work-related factors---not measured occupational stress---and is treated as potentially \emph{associated} with tinnitus burden and hearing-related outcomes rather than as a proven cause of cochlear injury. We report the \emph{primary} perceived-stress analysis in KNHANES (which carries a general stress item) together with a \emph{supporting} analysis in NHANES.

\textbf{Method (prespecified):} The Korea National Health and Nutrition Examination Survey (KNHANES) is the development dataset and the U.S. National Health and Nutrition Examination Survey (NHANES) is the geographic external-validation dataset; we fix the survey cycles, age windows, audiometric frequencies, exclusions, and expected analytic sample sizes in advance. A causal directed acyclic graph (DAG) defines a minimal sufficient adjustment set and separates confounders from mediators (e.g., depressed mood, sleep), with primary analyses avoiding overadjustment and full-adult versus employed-restricted analyses addressing the healthy-worker effect. Endpoints (tinnitus presence, bothersome tinnitus, better- and worse-ear pure-tone average, and asymmetry) are modeled with survey-weighted regression; two \emph{research-only} risk-stratification models are evaluated with discrimination, calibration-in-the-large and calibration slope (before and after recalibration), decision-curve utility, common-support checks, and subgroup fairness (including socioeconomic status and race/ethnicity), reported per STROBE and TRIPOD+AI. A deterministic red-flag layer for sudden sensorineural hearing loss (SSNHL) overrides model output; its evaluation prespecifies vignette composition, dual otolaryngology/audiology labeling, inter-rater agreement, and sensitivity with confidence intervals. An open medical language model (e.g., MedGemma) is confined to schema-constrained extraction verified by clinicians.

\textbf{Contributions:} (1) a survey-weighted public-data framework that distinguishes stress--tinnitus from stress--audiometric-threshold associations without unsupported causal claims; (2) a prespecified KNHANES$\rightarrow$NHANES external-validation plan reporting calibration, recalibration, clinical utility, and subgroup fairness; (3) a clinically governed AI architecture in which schema-constrained extraction is clinician-verified and a deterministic red-flag layer overrides probabilistic predictions for time-critical hearing symptoms; and (4) an open, reproducible scaffold of harmonized variable mappings, executable code, model cards, safety tests, and a synthetic-data pathway.

\textbf{Primary results (KNHANES, adults 40--69, cycles 2010--2012):} Survey-weighted prevalence was 21.1\% (95\% CI 19.9--22.3) for tinnitus and 15.4\% (14.5--16.4) for better-ear hearing loss ($>$25 dB HL), with high perceived stress at 24.3\% (23.2--25.3). In the mutually adjusted, minimal-sufficient model (perceived stress, occupational noise, age, sex), high perceived stress was associated with \emph{tinnitus} (OR 1.42, 1.26--1.60, $p<0.001$) but only weakly and non-significantly with \emph{audiometric hearing loss} (OR 1.18, 0.99--1.41, $p=0.07$), which was instead dominated by occupational noise (OR 1.72, 1.44--2.06) and age (OR 1.13/year). This symptom-versus-threshold dissociation is exactly the prespecified pattern of hypothesis H1. \textbf{Supporting results (NHANES):} using a PHQ-9 distress proxy (a mediator, since NHANES lacks a stress item), the same dissociation held. All estimates are associations, not causal effects.

\textbf{Conclusions:} Registered public-data analyses can rigorously characterize associations and benchmark explainable, safety-gated AI, but prospective clinical cohorts remain necessary for treatment timing, recovery, and prognosis. STARS fixes the design, safety governance, and reproducible scaffold, and demonstrates the primary perceived-stress finding end-to-end on real KNHANES data, with NHANES as cross-national support.
\end{abstract}

\noindent\textbf{Keywords:} tinnitus; bothersome tinnitus; hearing loss; audiometry; perceived stress; work-related factors; sudden sensorineural hearing loss; KNHANES; NHANES; external validation; explainable AI; open medical language models; TRIPOD+AI; study protocol; reproducibility; audiology.
